% HAFF: Holographic Alaya-Field Framework
% Postscript: participation, persistence and finite self-knowledge.
\documentclass[12pt,a4paper]{article}
\usepackage{amsmath,amssymb,amsfonts,amsthm}
\usepackage[hidelinks]{hyperref}
\usepackage{geometry}
\geometry{margin=1in}
\newtheorem{proposition}{Proposition}
\newtheorem{conjecture}{Conjecture}
\hypersetup{pdftitle={Postscript: Persistence, Participation and an Unfinished Self},pdfauthor={Sidong Liu}}
\title{Postscript: Persistence, Participation\\and an Unfinished Self}
\author{Sidong Liu, PhD\\iBioStratix Ltd\\\texttt{sidongliu@hotmail.com}}
\date{October 2026}
\begin{document}
\emergencystretch=2em
\raggedbottom
\maketitle

A world can leave a record without leaving anyone able to use it. A reader can continue without knowing everything about itself. These two distinctions frame the question of closure in the Holographic Alaya--Field Framework: what can end, what can remain, and what would be required for participation to begin again?

The framework studies physical access rather than treating intelligibility as an attribute of matter in isolation. A channel resolves some differences and loses others. A stable relation can preserve classical labels or quantum content. A body can use retained history to change its next action, its own channels and the physical sources it produces \cite{A,B,C,D,G,FoundationIII}. These claims connect world formation with agency, but they do not make the universe depend for its existence on a particular observer. Nor do they identify a stable observable algebra, a finite reconstruction of geometry and a spacetime metric. Each reaches a different part of the question.

\section{A trace is not a continuing reader}

Consider a fair classical label $X$ copied perfectly into $R$ two-level memories. Its joint state is
\[
 \rho_{XC}=\frac12\sum_{x=0}^1
 |x\rangle\langle x|\otimes
 (|x\rangle\langle x|)^{\otimes R}.
\]
Every memory alone is maximally mixed; the bank has no marginal purity. Nevertheless each copy determines $X$. Dephasing in the $Z$ basis leaves this joint record unchanged. The product state with the same maximally mixed marginals carries no such record. The difference is in the correlations, not in the local inventories.

An archive of this kind may outlast the machinery that formed it. A later body with appropriate access can still read it. Without that machinery, however, the archive does not acquire a new observation, test an interpretation or correct its own readout. Passive protection and active participation are different physical achievements. Foundation's distinction between held content and return into another operation makes this difference constructive \cite{Foundation}. The continuing operation must actually be available; a comparison state used to bound an error is not a free preparation of another program.

This is also a philosophical distinction. What persists as a trace need not persist as an agent. The past can remain physically relevant through future interventions even after its original reader is gone. Conversely, continuity of an agent need not mean preserving every past datum or every microscopic component. In the framework's agency account it means maintaining an organization that can carry usable relations into later choices, repair its channels and revise what those relations warrant \cite{C,TD}. This is a proposed criterion of causal participation, not a claim that causal organization alone explains experience.

\section{Acquisition and renewal have different limits}

HAFF~F bounds record acquisition at a declared initial resource boundary \cite{F}. Its assumptions include initially independent, unused memories and a specified thermal work account. The initial purity
\[
 B=\sum_k[\ln d_k-S(\sigma_k)]
\]
is then part of what can pay for new correlations, together with supplied resources and the system's own contribution. The example above explains why one cannot impose this product boundary anew at each time. Old correlations are resources even when $B=0$. Running out of fresh purity does not erase them.

Renewal asks another question: what must be done when an existing useful relation degrades? In F's conditional-independent bank, $C_i=X\oplus N_i$, with error probability $\varepsilon<1/2$, independent flips of probability $0<p\le1/2$ give
\[
 q=p+(1-2p)\varepsilon,\qquad
 d_{p,\varepsilon}=h(q)-h(\varepsilon)>0,
\]
where $h(u)=-u\ln u-(1-u)\ln(1-u)$. Returning the complete retained label, bank and controller to their joint target removes $R d_{p,\varepsilon}$ nats. An independent Gibbs bath at temperature $T$ must receive
\[
 Q\ge k_BT R d_{p,\varepsilon}.
\]
Here the flip record is unavailable, the bath is initially independent of the noisy retained state, and the controller cannot absorb the entropy without being restored. For $p=.12$ and $\varepsilon=.05$, eight copies of one fair label require $1.9025859874$ nats of correction. The number of distinct labels is one; the independent degradation of eight carriers still has a price.

A sole thermal chiral exhaust with the additional current assumptions in F and D has a quadratic finite-rate price. That result does not apply to every bath or every control pulse. Exact protection gives $d=0$ in the bank model and defeats a blanket assertion that records require positive upkeep. Accessible noise records, correlated errors and a controller allowed to consume its own purity alter the boundary as well. The distinction matters before any cosmological extrapolation.

\section{A finite-body extension}

The bank is not the only setting in which a retained relation has a renewal cost. The following result treats an entire finite memory as one object, allows correlated and quantum content, and makes the noise assumption explicit. Entropies and mutual informations are in nats.

\begin{proposition}[Renewing a useful relation under partial reset]
\label{prop:body}
Let $\theta_{XMC}$ be a target state of a retained reference $X$, a memory $M$ of dimension $d$, and a controller $C$. Suppose $I_\theta(M:X)\ge\iota>0$. During each degradation interval the unavailable noise acts as
\[
 \mathcal D_\lambda(\theta)
 =(1-\lambda)\theta+
 \lambda\,\theta_{XC}\otimes\tau_M,\qquad
 \tau_M=I_M/d,\quad 0<\lambda\le1.
\]
A correction restores the complete target $\theta_{XMC}$, using a bath initially independent of the noisy retained state and Gibbs at temperature $T$. If the correction is a possibly driven unitary on the retained system and bath, with no other entropy sink and vanishing endpoint interaction energy, then
\begin{equation}\label{eq:heatbody}
 Q\ge k_BT\lambda\iota.
\end{equation}
If the total retained-system energy, including internal couplings, is equal immediately before and after the correction, and driven work and the counted bath heat are the only energy flows, then $W=Q$. A degenerate memory with additive, unchanged apparatus energies is one sufficient realization. For $N$ renewals with $T_j\ge T_*>0$, $\lambda_j\ge\lambda_*>0$, and $I_{\theta_j}(M_j:X_j)\ge\iota_*>0$, a finite total supplied work $\sum_j W_j\le E_{\rm av}$ implies
\begin{equation}\label{eq:number}
 N\le\frac{E_{\rm av}}{k_BT_*\lambda_*\iota_*}.
\end{equation}
For the last statement each interval must satisfy the preceding cyclic correction and energy conditions, even when its target differs.
\end{proposition}

\begin{proof}
Write $Z=XMC$ and $\eta=\mathcal D_\lambda(\theta)$. Entropy concavity gives
\begin{align*}
 S(\eta)-S(\theta)
 &\ge \lambda[S(\theta_{XC})+\ln d-S(\theta)]\\
 &=\lambda[\ln d-S_\theta(M|XC)]\\
 &=\lambda[\ln d-S_\theta(M)+I_\theta(M:XC)]\\
 &\ge\lambda I_\theta(M:X)\ge\lambda\iota.
\end{align*}
The last line uses $S(M)\le\ln d$ and monotonicity of mutual information under discarding $C$. The proof is valid for quantum $X$ and for a controller correlated with the target.

Initially $Z$ and the Gibbs bath $B$ are independent. Entropy conservation under the correction and subadditivity at its end give $\Delta S_B\ge S(\eta)-S(\theta)$. For $\tau_B=e^{-\beta H_B}/\operatorname{Tr}e^{-\beta H_B}$,
\[
 \beta Q=\Delta S_B+D(\rho'_B\Vert\tau_B),
 \qquad \beta=(k_BT)^{-1}.
\]
This proves \eqref{eq:heatbody}. With vanishing endpoint bath coupling, the first law reads $W=\Delta E_{XMC}+Q$. The total-energy and energy-flow assumptions give $\Delta E_{XMC}=0$ and hence $W=Q$; summing the lower bounds proves \eqref{eq:number}.
\end{proof}

The reset mixture models a lost alternative, not an accessible branch flag the controller can undo for free. It acts on the whole designated memory, so no protected subsystem within that memory escapes the stated channel. The cost is paid for restoring the joint relation, not for restoring a local marginal. A fair classical copy has a maximally mixed memory marginal and still satisfies the bound. An initially independent maximally mixed memory has $\iota=0$ and does not.

This extends the renewal argument beyond independent copies. It does not establish that every physical body experiences this noise, that every update needs a fixed amount of information, or that every target can be restored under a body's allowed access. The bound is necessary where such a correction exists. A consumed controller, an accumulating waste register or an accessible auxiliary noise record must be included among the resources. If the noise environment is retained, its correlations cannot be discarded when defining the target. Calling only the memory cyclic would hide the expenditure. The entropy proof permits an externally scheduled unitary; it is not a proof of autonomous apparatus closure. A physical battery that takes part in the entropy account must be included or separately shown not to supply an additional entropy sink.

Equation~\eqref{eq:number} bounds successful renewals, not necessarily calendar time. A system can postpone a last update arbitrarily far. A finite upper bound on an uninterrupted active interval additionally requires a positive minimum update rate and a bound on gaps. Hibernation, cooling toward zero temperature, decreasing task quality and improving protection are substantive alternatives, not semantic exceptions.

\section{An unfinished self can continue}

A demand for complete self-description makes participation seem impossible for the wrong reason. A self-model need only support the inquiries and actions on which continuation depends. T-DOME~II gives a concrete example: history must retain the signed relation between the task and the body's channels. An accurately learned but unsigned axis can leave the prediction useless \cite{TD}. The useful self-frame therefore consists in how a system can reach its task, not in an inventory of everything it is.

Some distinctions are inaccessible to a body's legitimate references. Foundation~II separates that failure of identifiability from disagreement, and separates agreement from correctness \cite{Foundation}. A hidden symmetry can make one parameter unknowable while leaving a valuable intervention completely identifiable. In other cases inquiry resolves the ambiguity, but its preparation and maintenance consume resources. A continuing body must decide which error matters and which probe can repair it before the inquiry destroys the content it was meant to help. Brain Anchor's learning-versus-survival tradeoff supplies one physical version of this problem \cite{Brain}.

Finite self-knowledge is consequently not equivalent to paralysis. A body may track residual errors, revise its references and retain a margin of viable action without ever obtaining a complete account of itself. The stronger philosophical claim is that an imperfect, historically formed organization can be a continuing subject precisely through this revisability. Its identity lies partly in the maintained relation between history and possible action, not in an immutable perfect image.

A rival account could identify continuation with persistence of the stored pattern. It then misses a relevant difference between two copies of the same archive: one has working probes and can revise its routing after an error, while the other's channels have permanently failed. Their stored descriptions agree; their possible next histories do not. The active criterion explains that difference without demanding exact microscopic identity. Two functioning copies may both inherit a causal history, so this criterion does not by itself select a unique first-person successor. That further identity question, like experience, cannot be decided by the survival of a data pattern alone.

That claim still needs an account of first-person experience. Consciousness~III's explanatory gap cannot be crossed by renaming calibration as awareness \cite{Consciousness}. A candidate bridge must explain why these particular internal relations matter to experience, distinguish alternatives that make the same reports, and expose a relevant intervention. Neither Proposition~\ref{prop:body} nor a finite self-frame proves that a body feels anything. The thermodynamic argument addresses the support for continuation; the experience argument addresses what, if anything, that continuation is like.

\section{A cosmological participation-window conjecture}

The universe's duration and a body's active history need not have the same limit. Dyson's cooling and scaling scenario explicitly allows indefinitely continued life and communication from finite energy under its assumed laws \cite{Dyson}. Finite resources alone therefore do not establish a finite active window. In asymptotically de Sitter models a detector can encounter cosmological-horizon thermal radiation \cite{GH}. Krauss and Starkman examine how horizon restrictions, recoverable energy and memory requirements change the prospect of indefinite activity \cite{KS}. These are competing conditional mechanisms, not observational proof of the far future.

The framework suggests a more specific conjecture.

\begin{conjecture}[Resource-limited participation]
In a late-time causal domain with finite accessible usable work, a positive environmental temperature floor and unavoidable loss of task-relevant relations, any fixed-capacity body that continues at a fixed nonzero task quality can perform only finitely many successful inquiry--revision--action updates, even while passive traces and spacetime structure persist.
\end{conjecture}

The proposed mechanism is loss of correlations on which usable permissions depend. A retained history supports diagnosis; diagnosis changes routing or protection; the revised body supports another action and its physical source. Mixing damages the relations required to repeat that loop. Restoration exports entropy or consumes a fresh resource. A horizon may restrict replenishment and a thermal floor may prevent the cost per restoration from tending to zero. Exhaustion then ends this participation without erasing every trace.

The bank and partial-reset propositions establish this mechanism in two specified memory models. G supplies finite history-dependent source differences, while Imprint and Brain connect retained relations to subsequent inquiry and bodily change \cite{G,Imprint,Brain}. The conjecture reaches further: an interacting adaptive body can compress memories, migrate codes and replace its task. To prove it for a proposed body class one must obtain a uniform positive resource cost per successful update after optimizing those operations. Neither a count of bit erasures nor a total marginal purity inventory supplies that result automatically. The temperature floor, accessible work bound and irreducible relation loss must also be established for the actual cosmological domain. They are not consequences of selecting an algebra.

The conjecture is weakened or falsified by a body at the stipulated quality with an indefinitely usable protected relation, by vanishing renewal cost after admissible compression or migration, by a cooling law permitting a summable sequence of positive costs, or by unbounded replenishment. An accessible noise record can break the proposed loss mechanism. Discovery of one such regime would identify which part of the mechanism fails, rather than merely extending a numerical lifetime. A static surviving archive is not a counterexample to this active claim, but a continuing reader with the stated resources would be.

The sustained selection--response conjecture in D has different resource
quantifiers \cite{D}: it asks for families of longer runs whose installed
energy, size and operating domain may grow, while one body is reused
within each run. The participation conjecture here fixes capacity and
limits the total accessible work. Neither quantifier refutes the other.
Foundation~III's finite semiclassical mean-action source--response
construction \cite{FoundationIII} concerns a specified operating window;
it does not establish indefinitely renewed thermal operation.

\section{Origin, return and what closure would mean}

An unused record port and an exhausted one can be locally indistinguishable and have different futures. The former may be independent of the source and backed by a usable gradient. The latter can carry correlations with an old history or lack the resources needed for another operation. Allowed future interventions can distinguish them. Equal local states therefore do not identify origin with terminus.

The stronger origin proposal in G is physical formation of permissions: a body with locally adjustable relations encounters diagnostic tasks, retains their results and uses them to install a useful coupling \cite{G}. Exploration, a distinguishable value difference and enough survival time must be supplied by the primitive dynamics. This moves the problem beyond an external program name, without proving that every supplied executor arises in this way. Enlarging an executor introduces another joint interaction to explain; it does not prove a universal final boundary or forbid a finite system from carrying a useful description of itself.

A cosmological return would require more still. One concrete mechanism candidate is that contraction or another geometric transition creates newly accessible nonequilibrium gradients; interactions then allocate part of that work to fresh record ports and bodies, while old correlations and exported entropy remain accounted for in other degrees of freedom. A repeated active phase would be a renewal of resources and formation, not necessarily recurrence of the global state. The physical transition must have its own dynamics, show that the gradients can be used, and explain why old entropy does not obstruct the new phase. Local thermal swaps or recurrence of a finite register do not establish any of these cosmological steps. A bounce that restores density but no usable work would fail the mechanism; so would a transition requiring an uncharged deletion of the old history. No such return is forced by the participation-window conjecture.

Closure thus has several meanings that must not be confused. A record can become unreadable to one body; a body can lose the capacity to revise itself; a cosmic domain can cease to support new active readers. None entails the nonexistence of its remaining physical structures. Nor must an imperfect self finish knowing itself before it can live. The framework's most consequential prospect is that history changes what a body can do next, and that this relation can be maintained and revised for a time without being complete.

\bigskip
\noindent\textit{A history can remain after the means of continuing it have gone. A continuing self need not wait to become a finished account of itself.}

\begin{thebibliography}{99}
\small
\bibitem{A}
S.~Liu, \emph{Physical Access and Correlation Geometry: The Holographic Alaya--Field Framework}, (2026), series record DOI: 10.5281/zenodo.18361706.
\bibitem{B}
S.~Liu, \emph{Coupling, Records, and Stable Relations: Formation, Protected Content and the Lifetime of a Geometry}, (2026), series record DOI: 10.5281/zenodo.18367060.
\bibitem{C}
S.~Liu, \emph{Causation, Agency, and Existence: Historical Control and Maintained Intervention}, (2026), series record DOI: 10.5281/zenodo.18374805.
\bibitem{D}
S.~Liu, \emph{Selection, Maintenance and Gravitational Response: Modular Constraints and a Physical Return}, (2026), series record DOI: 10.5281/zenodo.18388881.
\bibitem{FoundationIII}
S.~Liu, \emph{Foundation III: The Perfuming Loop. Held Content, Learned Frames and Physical Return}, (2026), series record DOI: 10.5281/zenodo.19642189.
\bibitem{F}
S.~Liu, \emph{Records, Dissipation and Intervention: Conditions for a Temporal Direction}, (2026), series record DOI: 10.5281/zenodo.18400425.
\bibitem{G}
S.~Liu, \emph{Internal Selection and Revisable Organization: How History Changes Permissions and Physical Sources}, (2026), series record DOI: 10.5281/zenodo.18402907.
\bibitem{TD}
S.~Liu, \emph{T-DOME}, working trilogy (2026), series record DOIs: I, 10.5281/zenodo.18574342; II, 10.5281/zenodo.18579703; III, 10.5281/zenodo.18591771.
\bibitem{Foundation}
S.~Liu, \emph{Foundation}, three working papers (2026); series records: I (DOI: 10.5281/zenodo.19627786), II (DOI: 10.5281/zenodo.19641372), III (DOI: 10.5281/zenodo.19642189).
\bibitem{Consciousness}
S.~Liu, \emph{Finite Self-Models and the Explanatory Gap: What a Continuing Observer Can Know About Itself}, (2026), series record DOI: 10.5281/zenodo.19607724.
\bibitem{Imprint}
S.~Liu, \emph{Imprint}, four working papers (2026), series record DOIs: 10.5281/zenodo.20122718, 10.5281/zenodo.20151451, 10.5281/zenodo.20151508 and 10.5281/zenodo.20151594.
\bibitem{Brain}
S.~Liu, \emph{The Brain as a Plastic Observable-Algebra Network: Learning Shared Access from Sensorimotor History}, (2026), series record DOI: 10.5281/zenodo.20044329.
\bibitem{Dyson}
F.~J.~Dyson, \emph{Time without end: Physics and biology in an open universe}, Reviews of Modern Physics \textbf{51}, 447 (1979), DOI: \href{https://doi.org/10.1103/RevModPhys.51.447}{10.1103/RevModPhys.51.447}.
\bibitem{GH}
G.~W.~Gibbons and S.~W.~Hawking, \emph{Cosmological event horizons, thermodynamics, and particle creation}, Physical Review D \textbf{15}, 2738 (1977), DOI: \href{https://doi.org/10.1103/PhysRevD.15.2738}{10.1103/PhysRevD.15.2738}.
\bibitem{KS}
L.~M.~Krauss and G.~D.~Starkman, \emph{Life, The Universe, and Nothing: Life and Death in an Ever-Expanding Universe}, Astrophysical Journal \textbf{531}, 22--30 (2000), DOI: \href{https://doi.org/10.1086/308434}{10.1086/308434}; \href{https://arxiv.org/abs/astro-ph/9902189}{arXiv:astro-ph/9902189}.
\end{thebibliography}
\end{document}
